% Harville 22.1, 22.2.1, 22.4; Lang LA III and IV.
\section*{1. Linear maps, kernels, and images}

\notetarget{note:04:linear}All spaces in this unit are finite-dimensional real
linear spaces. A \textbf{linear map} \(T:V\to W\) satisfies
\[
T(ax+by)=aT(x)+bT(y)\qquad(x,y\in V,\ a,b\in\R).
\]
Here \(V\) is the domain and \(W\) the codomain. Taking \(a=b=0\)
gives \(T(0)=0\); repeated addition then gives
\[
T\left(\sum_{j=1}^k c_jv_j\right)=\sum_{j=1}^k c_jT(v_j),
\]
including the empty sum. The zero map \(x\mapsto0\) and identity map
\(I_V:x\mapsto x\) satisfy the defining equality directly.
For \(A:m\times n\), \(T_A:\R^n\to\R^m,\ x\mapsto Ax\), is linear by
the distributive laws. More generally, fixed \(A:m\times p,B:q\times n\) give
the linear map \(X\mapsto AXB\) from \(\R^{p\times q}\) to \(\R^{m\times n}\), since
\[
A(aX+bY)B=aAXB+bAYB.
\]

\notetarget{note:04:kernel-image}Define
\[
\ker T=\{x\in V:T(x)=0\},\qquad
\im T=T(V)=\{T(x):x\in V\}.
\]
The map is \textbf{injective} if \(T(x)=T(y)\) implies \(x=y\), and
\textbf{surjective} if \(\im T=W\). In the matrix case,
\(\ker T_A=N(A)\) and \(\im T_A=\mathcal C(A)\).

\begin{proposition}\label{prop:04:basic}\notetarget{note:04:basic}
For a linear map \(T:V\to W\):
\begin{enumerate}
\item \(\ker T\) and \(\im T\) are subspaces of \(V\) and \(W\).
\item \(T(\Span(v_1,\ldots,v_k))=\Span(T(v_1),\ldots,T(v_k))\).
\item \(T\) is injective if and only if \(\ker T=\{0\}\).
\end{enumerate}
\end{proposition}
\begin{proof}\notetarget{note:04:basic-proof}
\begin{itemize}
\item Both sets contain zero. If \(x,y\in\ker T\), then
  \(T(ax+by)=0\). If \(u=T(x),v=T(y)\), then
  \(au+bv=T(ax+by)\in\im T\). These are the subspace tests.
\item As \(c_1,\ldots,c_k\) range over \(\R\), the identity
  \[
  T\left(\sum_jc_jv_j\right)=\sum_jc_jT(v_j)
  \]
  gives every vector on each side of (b). Empty lists give \(\{0\}\).
\item An injective map has only one preimage of zero, namely zero.
  Conversely, \(T(x)=T(y)\) gives \(x-y\in\ker T\);
  if that kernel is \(\{0\}\), then \(x=y\).
\end{itemize}
\end{proof}

\begin{theorem}[Prescribing the images of a basis]\label{thm:04:prescribed}
\notetarget{note:04:prescribed}
Let \(B=(v_1,\ldots,v_n)\) be a basis of \(V\) and let
\(w_1,\ldots,w_n\) be any vectors of \(W\).
There is exactly one linear map \(T:V\to W\) such that \(T(v_j)=w_j\).
\end{theorem}
\begin{proof}\notetarget{note:04:prescribed-proof}
\begin{itemize}
\item Each \(v\in V\) has unique coefficients \(v=\sum_jx_jv_j\).
  Define \(T(v)=\sum_jx_jw_j\). Uniqueness of the coefficients makes
  this a well-defined map.
\item If \(v=\sum_jx_jv_j,\ u=\sum_jy_jv_j\), the coefficients of
  \(av+bu\) are \(ax_j+by_j\). Hence
  \[
  T(av+bu)=\sum_j(ax_j+by_j)w_j=aT(v)+bT(u).
  \]
  In particular \(T(v_j)=w_j\).
\item Any linear map with these basis images must send
  \(\sum_jx_jv_j\) to \(\sum_jx_jw_j\), so it equals \(T\).
\end{itemize}
For \(V=\{0\}\), the empty basis prescribes the unique map \(0\mapsto0\).
\end{proof}

\notetarget{note:04:basis-images}Thus two linear maps agreeing on a basis
agree everywhere. More generally, agreement on a spanning list suffices:
expand any vector in that list and apply linearity to each summand.

\begin{theorem}[Rank--nullity]\label{thm:04:nullity}\notetarget{note:04:nullity}
For every linear map \(T:V\to W\),
\[
\dim V=\dim\ker T+\dim\im T.
\]
The number \(\dim\im T\) is the \textbf{rank} of \(T\), written \(\rank T\).
\end{theorem}
\begin{proof}\notetarget{note:04:nullity-proof}
\begin{itemize}
\item Extend a basis \((z_1,\ldots,z_s)\) of \(\ker T\) to a basis
  \((x_1,\ldots,x_r,z_1,\ldots,z_s)\) of \(V\), using Unit 3,
  Theorem~\ref{three-thm:03:extract}.
\item Write any \(v\in V\) as
  \(v=\sum_i a_ix_i+\sum_jb_jz_j\). Then
  \(T(v)=\sum_i a_iT(x_i)\); hence the \(T(x_i)\) span \(\im T\).
\item Suppose \(\sum_i c_iT(x_i)=0\). Then
  \(\sum_i c_ix_i\in\ker T\), so \(\sum_i c_ix_i=\sum_jd_jz_j\).
  Independence of the extended basis forces every \(c_i,d_j\) to be zero.
  Thus \(T(x_1),\ldots,T(x_r)\) are independent.
\item They form an image basis. Consequently
  \(\dim\im T=r,\ \dim\ker T=s,\ \dim V=r+s\).
\end{itemize}
The argument includes either empty list and the zero domain.
\end{proof}

\notetarget{note:04:isomorphism}A linear bijection is an \textbf{isomorphism}.
\begin{itemize}
\item Its inverse is linear: if \(y=T(x),y'=T(x')\), then
\[
T^{-1}(ay+by')=T^{-1}(T(ax+bx'))=ax+bx'.
\]
\item If \((v_j)\) is a basis of \(V\), surjectivity gives
\(w=T(v)=\sum_jc_jT(v_j)\) for every \(w\in W\). Thus the images span \(W\).
For independence,
\[
\sum_jc_jT(v_j)=0
\ \Longrightarrow\ T\left(\sum_jc_jv_j\right)=0
\ \Longrightarrow\ \sum_jc_jv_j=0
\ \Longrightarrow\ c_j=0\quad(\forall j).
\]
Here the last two implications use injectivity and the original basis.
An isomorphism therefore sends a basis to a basis, so \(\dim V=\dim W\).
\item Conversely, given bases \((v_j)\), \((w_j)\) of the same length,
Theorem~\ref{thm:04:prescribed} constructs \(T(v_j)=w_j\).
Every \(\sum_jc_jw_j\) has the unique preimage \(\sum_jc_jv_j\),
so \(T\) is an isomorphism.
\end{itemize}

\begin{proposition}\label{prop:04:equal-dim}\notetarget{note:04:equal-dim}
If \(\dim V=\dim W\), a linear map \(T:V\to W\) is injective
if and only if it is surjective.
\end{proposition}
\begin{proof}\notetarget{note:04:equal-dim-proof}
By Proposition~\ref{prop:04:basic} and rank--nullity,
\[
T\text{ injective}
\Longleftrightarrow\dim\ker T=0
\Longleftrightarrow\dim\im T=\dim V=\dim W.
\]
Since \(\im T\subseteq W\), the last equality is equivalent to
\(\im T=W\), by the equal-dimension subspace criterion of Unit 3.
\end{proof}

% Lang LA III §3.15(a,b,c), Shakarchi III §3.15; image is an added application.
\begin{example}[Symmetrization]\label{example:04:symmetrization}
\notetarget{note:04:symmetrization}
On \(\R^{n\times n}\), \(n\geqslant1\), define
\[
P(A)=\frac{A+A^\top}{2}.
\]
(a) Prove linearity. (b) Find \(\ker P\). (c) Find its dimension.
Also determine \(\im P\) and \(\rank P\).
\par\smallskip
\notetarget{note:04:symmetrization-solution}
(a) Transpose is linear, so
\[
P(aA+bB)=\frac{aA+bB+aA^\top+bB^\top}{2}=aP(A)+bP(B).
\]
(b) \(P(A)=0\) is exactly \(A^\top=-A\): the kernel is the skew-symmetric space.
(c) Such matrices have zero diagonal and \(a_{ji}=-a_{ij}\). Thus
\[
A=\sum_{i<j}a_{ij}(E_{ij}-E_{ji}).
\]
These matrices are independent by their entries above the diagonal.
They form a basis with \(n(n-1)/2\) members: the \(n(n-1)\) ordered
distinct pairs occur in reversal pairs.
Every \(P(A)\) is symmetric; conversely, a symmetric \(S\) satisfies
\(P(S)=S\). Therefore the image is the symmetric space and
\[
\rank P=n^2-\frac{n(n-1)}2=\frac{n(n+1)}2.
\]
\end{example}

\section*{2. Coordinates and matrix representations}

\notetarget{note:04:coordinates}For an \textbf{ordered basis}
\(B=(v_1,\ldots,v_n)\), the coordinate column \([v]_B\) is defined by
\[
v=\sum_jx_jv_j,\qquad [v]_B=(x_1,\ldots,x_n)^\top.
\]
The coordinate map \(v\mapsto[v]_B\) is a linear bijection onto \(\R^n\).
Indeed, coefficients of \(av+bu\) are \(a[v]_B+b[u]_B\);
equal coordinate columns give equal vectors, and every \(x\in\R^n\)
is the coordinate column of \(\sum_jx_jv_j\).
The order is part of the basis data.
For the zero space, \(\R^0\) contains just the empty coordinate column.
The unique \(0\times0\) matrix is \(I_0\); its product with itself is
\(I_0\), so \(I_0^{-1}=I_0\). Products with empty inner index sets use
the empty-sum convention of Unit 3.

\begin{theorem}[Matrix of a linear map]\label{thm:04:representation}
\notetarget{note:04:representation}
Let \(B=(v_1,\ldots,v_n)\) and \(C=(w_1,\ldots,w_m)\) be ordered bases
of \(V,W\). For a linear map \(T:V\to W\), put
\[
A=[T]_{C\leftarrow B}
 =\bigl[\,[T(v_1)]_C\ \cdots\ [T(v_n)]_C\,\bigr]\in\R^{m\times n}.
\]
Then \([T(v)]_C=A[v]_B\) for every \(v\in V\).
This identity determines \(A\) uniquely.
Conversely, every \(m\times n\) matrix is the matrix of exactly one
linear map in these bases.
\end{theorem}
\begin{proof}\notetarget{note:04:representation-proof}
\begin{itemize}
\item Write \(T(v_j)=\sum_i a_{ij}w_i\) and \(v=\sum_jx_jv_j\). Then
  \[
  T(v)=\sum_jx_j\sum_i a_{ij}w_i
  =\sum_i\left(\sum_j a_{ij}x_j\right)w_i.
  \]
  Uniqueness of the \(C\)-coefficients gives \([T(v)]_C=Ax\).
\item Applying the identity to \(v_j\), whose coordinate column is \(e_j\),
  determines each column of \(A\).
\item Given \(A\), prescribe \(T(v_j)=\sum_i a_{ij}w_i\).
  Theorem~\ref{thm:04:prescribed} gives existence and uniqueness.
\end{itemize}
\end{proof}

\notetarget{note:04:rank-coordinates}Coordinates also identify the two
kernels and the two images:
\[
\{[v]_B:v\in\ker T\}=N(A),\qquad
\{[w]_C:w\in\im T\}=\mathcal C(A).
\]
For the first equality, \(T(v)=0\) is equivalent to \(A[v]_B=0\),
and every coordinate column has a preimage \(v\).
For the second, each image coordinate is \(A[v]_B\); conversely,
any \(Ax\) is the coordinate column of \(T(\sum_jx_jv_j)\).
The restricted coordinate maps are linear bijections. Since such maps
preserve bases, they preserve dimension, giving
\[
\rank T=\rank A,\qquad \dim\ker T=\dim N(A).
\]

% Lang LA IV §3 Example 1; original numerical values retained.
\begin{example}\label{example:04:representation}\notetarget{note:04:matrix-example}
Suppose \(B=(v_1,v_2)\), \(C=(w_1,w_2,w_3)\) are ordered bases
of \(V,W\), and a linear map \(T:V\to W\) satisfies
\[
T(v_1)=3w_1-w_2+17w_3,\qquad T(v_2)=w_1+w_2-w_3.
\]
Find the representing matrix and the image of \(x_1v_1+x_2v_2\).
\par\smallskip
\notetarget{note:04:matrix-example-solution}
The coefficients of the two displayed images are the two columns:
\[
[T]_{C\leftarrow B}=
\begin{pmatrix}3&1\\-1&1\\17&-1\end{pmatrix}.
\]
Consequently
\[
T(x_1v_1+x_2v_2)
=(3x_1+x_2)w_1+(-x_1+x_2)w_2+(17x_1-x_2)w_3.
\]
\end{example}

\notetarget{note:04:map-operations}For maps \(T,S:V\to W\), define
\((aT+bS)(v)=aT(v)+bS(v)\). If both maps are linear, expansion at
\(\alpha u+\beta v\) gives
\[
(aT+bS)(\alpha u+\beta v)
=\alpha(aT+bS)(u)+\beta(aT+bS)(v).
\]
Thus their linear combination is linear; its values on basis members give
\([aT+bS]_{C\leftarrow B}=a[T]_{C\leftarrow B}+b[S]_{C\leftarrow B}\).

\begin{proposition}[Composition and inverses]\label{prop:04:composition}
\notetarget{note:04:composition}
Let \(T:V\to W\) and \(S:W\to U\) be linear, with representing matrices
\(A=[T]_{C\leftarrow B}\) and \(D=[S]_{E\leftarrow C}\).
The composition \(ST:v\mapsto S(T(v))\) is linear, and
\[
[ST]_{E\leftarrow B}=DA.
\]
Also \([I_V]_{B\leftarrow B}=I_n\). If \(T\) is an isomorphism,
\[
[T^{-1}]_{B\leftarrow C}=A^{-1}.
\]
\end{proposition}
\begin{proof}\notetarget{note:04:composition-proof}
\begin{itemize}
\item Successive linearity gives
  \(ST(av+bu)=S(aT(v)+bT(u))=aST(v)+bST(u)\).
\item For every \(v\),
  \([ST(v)]_E=D[T(v)]_C=DA[v]_B\).
  Uniqueness of the representing matrix gives the formula.
\item The identity sends \(v_j\) to \(v_j\), giving the columns \(e_j\).
  The inverse map is linear, as already proved.
  Its matrix \(H\) satisfies \(HA=I_n\) and \(AH=I_m\) from
  \(T^{-1}T=I_V,\ TT^{-1}=I_W\).
  An isomorphism preserves dimension, so \(m=n\), and \(H=A^{-1}\).
\end{itemize}
\end{proof}

\begin{theorem}[Changing bases]\label{thm:04:change}\notetarget{note:04:change}
Keep \(A=[T]_{C\leftarrow B}\). Let \(E=(u_1,\ldots,u_n)\)
and \(F=(z_1,\ldots,z_m)\) be new bases of \(V,W\). Set
\[
R=[I_V]_{B\leftarrow E},\qquad Q=[I_W]_{C\leftarrow F}.
\]
Then \(R,Q\) are invertible, and
\[
[v]_B=R[v]_E,\qquad [w]_C=Q[w]_F,\qquad
[T]_{F\leftarrow E}=Q^{-1}AR.
\]
\end{theorem}
\begin{proof}\notetarget{note:04:change-proof}
\begin{itemize}
\item The first two identities are the representation formula applied
  to identity maps. Composing the coordinate changes in reverse order
  gives \(I_n\) or \(I_m\); hence both changes are invertible.
\item For every \(v\in V\),
  \[
  Q[T(v)]_F=[T(v)]_C=A[v]_B=AR[v]_E.
  \]
  Multiplying by \(Q^{-1}\) gives the claimed matrix.
\end{itemize}
\end{proof}

\notetarget{note:04:similarity}For \(T:V\to V\), using the same basis
on domain and codomain gives the special rule
\[
[T]_{E\leftarrow E}=R^{-1}[T]_{B\leftarrow B}R.
\]
Matrices \(A,H\) are called \textbf{similar} if \(H=R^{-1}AR\) for an
invertible \(R\).
Every such \(R\) is a change-of-basis matrix: define \(u_j\) by
\([u_j]_B=Re_j\). Its columns form a basis of \(\R^n\), and the inverse
coordinate map sends that basis to \((u_j)\).
Thus similar matrices represent the same linear map in different bases.
In particular they have equal rank, by the coordinate rank equality above.

% Lang LA IV §3.1(a); manual gives final matrix only, intermediate steps supplied.
\begin{example}[A change of coordinates]\label{example:04:change}
\notetarget{note:04:change-example}
In \(\R^3\), take the ordered lists
\[
B=((1,1,0)^\top,(-1,1,1)^\top,(0,1,2)^\top),\qquad
C=((2,1,1)^\top,(0,0,1)^\top,(-1,1,1)^\top).
\]
Verify that they are bases and find \(P=[I]_{C\leftarrow B}\).
\par\smallskip
\notetarget{note:04:change-example-solution}
Let \(U,V\) contain the columns of \(B,C\):
\[
U=\begin{pmatrix}1&-1&0\\1&1&1\\0&1&2\end{pmatrix},
\qquad V=\begin{pmatrix}2&0&-1\\1&0&1\\1&1&1\end{pmatrix}.
\]
If \(U(a,b,c)^\top=0\), its first and third equations give \(a=b,\ b=-2c\);
the second gives \(a+b+c=-3c=0\). Thus all three coefficients are zero,
and the columns of \(U\) are a basis.
For any \(u=(u_1,u_2,u_3)^\top\), solve \(Vp=u\):
\[
2p_1-p_3=u_1,\qquad p_1+p_3=u_2,\qquad p_1+p_2+p_3=u_3.
\]
These have the unique solution
\[
p_1=(u_1+u_2)/3,\qquad p_2=u_3-u_2,\qquad
p_3=(2u_2-u_1)/3,
\]
also proving that \(C\) is a basis. Substitute the three columns of \(U\):
\[
P=\begin{pmatrix}2/3&0&1/3\\-1&0&1\\1/3&1&2/3\end{pmatrix},
\qquad VP=U.
\]
The direction is fixed by \([v]_C=P[v]_B\).
\end{example}

\section*{3. Restrictions, product ranks, and matrix equations}

\notetarget{note:04:restriction}For a subspace \(U\subseteq V\),
the \textbf{restriction} \(T|_U:U\to W\) is defined by
\(T|_U(x)=T(x)\) for \(x\in U\).
It is linear: \(ax+by\in U\) and
\(T|_U(ax+by)=aT|_U(x)+bT|_U(y)\).
Its kernel and image are
\[
\ker(T|_U)=\{x\in U:T(x)=0\}=U\cap\ker T,
\qquad \im(T|_U)=\{T(x):x\in U\}=T(U).
\]

\begin{theorem}[Rank of a composition]\label{thm:04:restriction}
\notetarget{note:04:product-rank}
For linear maps \(S:V\to W,\ T:W\to Z\),
\[
\rank(TS)=\rank S-\dim(\im S\cap\ker T).
\]
In particular, for \(A:m\times n,\ B:n\times p\),
\[
\rank(AB)=\rank B-\dim(\mathcal C(B)\cap N(A)).
\]
\end{theorem}
\begin{proof}\notetarget{note:04:product-rank-proof}
\begin{itemize}
\item Put \(U=\im S\). Then
\(\ker(T|_U)=U\cap\ker T\) and
\(\im(T|_U)=T(S(V))=\im(TS)\).
\item Rank--nullity for \(T|_U\) gives
\[
\rank(TS)=\dim U-\dim(U\cap\ker T)
=\rank S-\dim(\im S\cap\ker T).
\]
\item Take \(S(x)=Bx,\ T(y)=Ay\) to obtain the matrix formula.
\end{itemize}
\end{proof}

\begin{proposition}[Product-rank bounds]\label{prop:04:bounds}
\notetarget{note:04:bounds}
For \(A:m\times n,\ B:n\times p\),
\[
\rank A+\rank B-n
\leqslant\rank(AB)\leqslant\min(\rank A,\rank B).
\]
Moreover,
\[
AB=0\Longleftrightarrow\mathcal C(B)\subseteq N(A);
\quad AB=0\Longrightarrow\rank A+\rank B\leqslant n.
\]
\end{proposition}
\begin{proof}\notetarget{note:04:bounds-proof}
\begin{itemize}
\item The intersection in Theorem~\ref{thm:04:restriction} has dimension
  between \(0\) and \(\dim N(A)=n-\rank A\). Substitution gives the
  lower bound and the upper bound \(\rank(AB)\leqslant\rank B\).
\item Every column combination \(ABx=A(Bx)\) lies in \(\mathcal C(A)\).
  Hence \(\mathcal C(AB)\subseteq\mathcal C(A)\), giving the other upper bound.
\item \(AB=0\) means that \(A\) annihilates every column of \(B\).
  By linearity this is equivalent to annihilating their entire span.
  The final inequality follows by setting \(\rank(AB)=0\) in the lower bound.
\end{itemize}
\end{proof}

\notetarget{note:04:rank-equality}The exact formula also shows
\[
\rank(AB)=\rank B
\Longleftrightarrow\mathcal C(B)\cap N(A)=\{0\}.
\]
Indeed, equality holds exactly when the subtracted dimension is zero;
a subspace has dimension zero exactly when it is \(\{0\}\).
The lower bound is called \textbf{Sylvester's inequality}.

% MA 4.16; original counterexample retained.
\begin{example}\label{example:04:ab-ba}\notetarget{note:04:ab-ba}
Must \(AB\) and \(BA\) have the same rank for square matrices \(A,B\)?
Consider
\[
A=\begin{pmatrix}1&0\\1&0\end{pmatrix},
\qquad B=\begin{pmatrix}0&0\\1&1\end{pmatrix}.
\]
\par\smallskip
\notetarget{note:04:ab-ba-solution}
Direct multiplication gives
\[
AB=0,\qquad BA=\begin{pmatrix}0&0\\2&0\end{pmatrix}.
\]
Their ranks are \(0\) and \(1\): the second matrix has one nonzero column.
Thus their ranks can differ, although Unit 2 gives
\(\tr(AB)=\tr(BA)\).
\end{example}

\begin{proposition}[Frobenius's inequality]\label{prop:04:frobenius}
\notetarget{note:04:frobenius}
For \(A:m\times n,\ B:n\times p,\ C:p\times q\),
\[
\rank(AB)+\rank(BC)\leqslant\rank B+\rank(ABC).
\]
\end{proposition}
\begin{proof}\notetarget{note:04:frobenius-proof}
\begin{itemize}
\item Set \(U=\mathcal C(B)\), \(U'=\mathcal C(BC)\).
Since \(BCx=B(Cx)\), we have \(U'\subseteq U\).
\item Apply Theorem~\ref{thm:04:restriction} twice:
\[
\begin{aligned}
\rank B-\rank(AB)&=\dim(U\cap N(A)),\\
\rank(BC)-\rank(ABC)&=\dim(U'\cap N(A)).
\end{aligned}
\]
\item The inclusion \(U'\cap N(A)\subseteq U\cap N(A)\) gives
\[
\rank(BC)-\rank(ABC)\leqslant\rank B-\rank(AB).
\]
Rearrange to obtain the stated inequality.
\end{itemize}
\end{proof}

\begin{theorem}[A matrix of right-hand sides]\label{thm:04:matrix-equation}
\notetarget{note:04:matrix-equation}
Let \(A:m\times n\), \(B:m\times q\), with \(m,n,q\geqslant1\),
and let \(X:n\times q\) be unknown.
Then
\[
AX=B\text{ is consistent}
\Longleftrightarrow\mathcal C(B)\subseteq\mathcal C(A)
\Longleftrightarrow\rank[\,A\ B\,]=\rank A.
\]
If \(X_0\) is a solution and the columns of \(K:n\times(n-r)\)
form a basis of \(N(A)\), where \(r=\rank A\), all solutions are
\[
X=X_0+KZ,\qquad Z\in\R^{(n-r)\times q}.
\]
Different \(Z\) give different solutions. In particular, a consistent
equation has a unique solution exactly when \(\rank A=n\).
\end{theorem}
\begin{proof}\notetarget{note:04:matrix-equation-proof}
\begin{itemize}
\item Writing \(X=[\,x_1\ \cdots\ x_q\,]\), \(B=[\,b_1\ \cdots\ b_q\,]\)
  gives \(AX=B\) exactly when \(Ax_j=b_j\) for every \(j\).
  This is equivalent to all columns of \(B\) lying in \(\mathcal C(A)\).
\item Since \(\mathcal C(A)\subseteq\mathcal C([\,A\ B\,])\), these
  spaces are equal exactly when their dimensions agree.
  Equality is equivalent to \(\mathcal C(B)\subseteq\mathcal C(A)\).
\item \(AK=0\), so \(A(X_0+KZ)=B\) for every \(Z\).
  Conversely, \(AX=B=AX_0\) gives \(A(X-X_0)=0\).
  Expand each column of \(X-X_0\) in the basis columns of \(K\);
  collecting these coefficients gives \(X-X_0=KZ\).
\item Independence of the columns of \(K\) gives
  \(KZ=KZ'\Rightarrow Z=Z'\), by comparing columns.
  There are \((n-r)q\) independent scalar parameters.
  For \(r=n\), the empty product \(KZ\) is zero.
  For \(r<n\) and \(q\geqslant1\), changing a parameter gives a different solution.
\end{itemize}
\end{proof}

% Bapat 2.16: original dimensions; the answer's counting hint is expanded.
\begin{example}\label{example:04:matrix-kernel}\notetarget{note:04:matrix-kernel}
Given \(A:9\times7\) and \(B:4\times3\), prove that there is a nonzero
\(7\times4\) matrix \(X\) satisfying \(AXB=0\).
\par\smallskip
\notetarget{note:04:matrix-kernel-solution}
The map
\[
\Phi:\R^{7\times4}\longrightarrow\R^{9\times3},\qquad X\longmapsto AXB
\]
is linear by the distributive calculation at the start of the unit.
The matrix-unit bases give dimensions \(28\) and \(27\). Since
\(\im\Phi\) is a subspace of the codomain, \(\rank\Phi\leqslant27\).
Consequently
\[
\dim\ker\Phi=28-\rank\Phi\geqslant1.
\]
The kernel therefore contains a nonzero matrix, and every such matrix
satisfies \(AXB=0\).
\end{example}
